\documentclass[11pt,letterpaper]{article}
\usepackage[T1]{fontenc}
\usepackage{lmodern}
\usepackage[margin=1.05in]{geometry}
\usepackage{amsmath,amssymb,amsthm,booktabs,microtype,xurl}
\usepackage[hidelinks,pdftitle={Nonhalting and PA-independence for a 209-state Turing machine},
  pdfauthor={James Chang}]{hyperref}

\newtheorem{theorem}{Theorem}[section]
\newtheorem{lemma}[theorem]{Lemma}
\newtheorem{proposition}[theorem]{Proposition}
\newcommand{\PA}{\mathrm{PA}}
\newcommand{\N}{\mathbb{N}}
\newcommand{\Con}{\operatorname{Con}}
\newcommand{\Prov}{\operatorname{Prov}}
\newcommand{\pair}[2]{\langle #1,#2\rangle}
\newcommand{\ax}{\mathsf{AX}}
\newcommand{\mpc}{\mathsf{MP}}
\newcommand{\shift}{\mathsf{S}}
\newcommand{\push}{\mathsf{PUSH}}
\newcommand{\pophead}{\mathsf{POP}}
\newcommand{\discard}{\mathsf{DROP}}
\newcommand{\src}[2]{\href{https://github.com/jamesyc/BusyBeaver/blob/main/#1}{\texttt{\nolinkurl{#2}}}}
\setlength{\emergencystretch}{2em}

\title{Nonhalting and PA-independence\\for a 209-state Turing machine}
\author{James Chang}
\date{}

\begin{document}
\maketitle

\begin{abstract}
We give an explicit two-symbol Turing machine with 209 working states
that searches for a contradiction in ordinary first-order Peano
arithmetic. The machine runs forever on a blank tape, and Peano
arithmetic proves neither its exact arithmetic halting sentence nor
its negation. The construction connects a register-program simulation
to a refutation compiler that works in every model of Peano arithmetic.
The literal transition table, both unprovability directions, and
blank-tape nonhalting are formalized in Lean~4.
\end{abstract}

\section{The result}\label{sec:result}

Busy Beaver problems ask how long a halting machine of a fixed size can
run~\cite{rado}. Explicit small machines whose behavior is independent of
a specified theory also give concrete examples of incompleteness.
Yedidia and Aaronson constructed such a machine for ZFC~\cite{yedidia}.
Here the theory is ordinary first-order Peano arithmetic with full
induction, denoted by $\PA$. This construction reduces the working-state
count from the earlier 211-state construction for $\PA$~\cite{pa211}.
The register engine draws on the K230 construction~\cite{k230}, and
its proof calculus uses ingredients from the PA238 development~\cite{pa238}
and Reshetnikov's restricted implication proof~\cite{proveit}.

Write $M$ for the machine in Section~\ref{sec:machine}. Let $H$ be the
sentence of arithmetic defined in Section~\ref{sec:arith}, and put
$N=\neg H$. Write $\PA\vdash A$ for the existence of a finite formal
derivation of $A$ from $\PA$. The statement $\Con(\PA)$, used inside
arithmetic, is the syntactic consistency sentence for the standard
provability predicate.

\begin{theorem}\label{thm:main}
The machine $M$ runs forever from the blank tape. Its arithmetic
nonhalting sentence is true in $\N$, and
\[
                 \PA\nvdash H \qquad\text{and}\qquad \PA\nvdash N .
\]
\end{theorem}

Under the external assumption of $\PA$ consistency, the machine runs
forever and $\PA\nvdash N$. The proof of $\PA\nvdash H$ uses the
soundness of $\PA$ in the standard natural numbers. In ordinary classical
mathematics that standard model also establishes consistency. Lean
uses the standard-model theorems for Foundation's theory
\texttt{Peano}~\cite{foundation}. The published statement has an explicit
presentation of the ordinary $\PA$ proof calculus, with its equivalence
to Foundation's presentation proved in the Solution.

\section{The machine and its startup}\label{sec:machine}

The tape is infinite in both directions, with alphabet $\{0,1\}$ and
initially all cells zero. State $0$ is the initial state. The working
state labels are
\[
            Q=\{0,\ldots,213\}\setminus\{50,122,150,157,205\}.
\]
A row $(q,b,w,D,q')$ reads $b$ at state $q$, writes $w$, moves one cell in
direction $D$, and enters $q'$. There is one row for every pair
in $Q\times\{0,1\}$. The state $215$ terminates the run. The 418 rows are
given in \src{data/pa209.tm}{data/pa209.tm}, with SHA-256 digest
\begin{quote}\small
\nolinkurl{0610d978675320f0b4b1ef34619dcd3daa5912632d7037e200dfcc70954fea51}.
\end{quote}
For example, the first row is $(0,0,1,R,1)$ and the transition
$(68,1,1,R,215)$ enters the halt state. Lean checks closure of all row
targets in $Q\cup\{215\}$ and the count of distinct working states.

The exact startup takes $58{,}210$ tape steps. It reaches the frame for
source site $25$ at literal control $81$, with the candidate register
equal to $2$ and every other logical register zero. The complete
register file, in its physical order, is
\[
              (0,0,0,0,0,0,2,0,0,1^{[26]}),
\]
where $1^{[26]}$ denotes 26 further entries equal to $1$. The prefix
checks the full tape and head, including those spare entries. The
definitions and proof are in \src{M6/Startup/Data.lean}{M6/Startup/Data.lean}
and \src{M6/Startup.lean}{M6/Startup.lean}.

\section{Simulation of the register program}\label{sec:physical}

The 3,948-site physical source is obtained by flattening the binary tree
in \src{M6/Program.lean}{M6/Program.lean}. Its instructions increment a
register, decrement a register or branch on zero, transfer control to
a structured loop boundary, or halt. The source crosswalk collapses the
extra physical site and normalizes the control transfers to the
3,947-site command program. The checked clean frame corresponds to the
zero command store at candidate $2$ in that program.

A frame contains the source path, header and marker, and a register
file in which a value $v$ is represented by a block $0\,1^{v+1}$.
The physical proof relabels the literal controls, sending the start
and halt to $74$ and $254$ respectively. The relabeling commutes with
each step, preserving the two tape lists. The resulting contracts
are proved for every source site and for both outcomes of a conditional
decrement. Each contract takes a frame to the complete frame of its
source successor in a positive number of tape steps, or implements the
source halt. The proof combines the finite control diagrams with the
unary arithmetic and register-file lemmas. It checks source addresses,
control destinations, and register effects separately.

The two right moves $(181,1,1,R,13)$ and $(187,0,0,R,75)$ supply the
changed copy operations used by the sixth axiom head in
Section~\ref{sec:commands}. The second move reaches control $75$ with
the preceding bits $10$. The following are its two preserving control
paths, with offsets relative to the head on entry at $75$.
\[
\begin{array}{c@{\quad}l}
0 & 75@0,\ 143@{-1},\ 149@{-2},\ 113@{-1},\ 146@0,\ 82@1,\\
1 & 75@0,\ 139@{-1},\ 171@0,\ 83@1.
\end{array}
\]
The resulting pair of providers $(82,83)$ performs the composition
whose first branch copies the register holding the second axiom
argument into the arithmetic accumulator. The branch entering $13$
uses the same copy operation. The four changed source operands are
decrement, increment, decrement, and increment, replacing register
$3$ by register $5$ within the two copies. The checked contracts verify
the actual entry histories, branch exits, and intervening providers
against the complete table. The source crosswalk and the source
command specification check the resulting logical argument order.

\begin{lemma}[Simulation]\label{lem:simulation}
If the command program from the clean candidate-$2$ store does not
halt, then the literal machine does not halt from the blank tape.
\end{lemma}

\begin{proof}
Append the certified startup to the positive-length simulation of a
source prefix. A source execution with arbitrarily long prefixes gives
tape runs of arbitrarily large length from the same initial
configuration. Determinism excludes a finite terminal run. The proof
uses the all-site coverage theorem and the crosswalk to the command
program.
\end{proof}

The model-valued version has an additional role. It describes the same
tape operations by arithmetic relations and concatenates their coded
traces in a model of $\PA$. The certificate proof in
Section~\ref{sec:refutation} uses that version.

\section{Proof search and soundness}\label{sec:commands}

\subsection{Proof commands}

Let $p(a,b)=\binom{a+b+1}{2}+b$ be Cantor pairing. Formula codes use
$I(a,b)=4p(a,b)$ for implication and $U(a,b)=4p(a,b)+2$ for
universal quantification. The two atomic forms have residues $1$ and
$3$ for membership and equality. Code $0$ is truth, and code $1$
decodes as $x_0\in x_0$. The decoder interprets implication into
code $1$ as negation. The calculus uses modus ponens and
generalization, and its axiom formulas are the exact constructors in
\src{M6/Source/ActualAXCode.lean}{ActualAXCode.lean}. The eighth head
requires a dominance condition on its fourth parameter. The
separation formula at the gate has the two indicated freshness
conditions on the third variable.

A store $(t,w,d,h)$ has two visible formula codes $t,h$, a stack $w$
encoded by iterated pairing, and a variable-label offset $d$.
Write $p_L,p_R$ for the total Cantor projections. The three
commands are as follows.
\begin{center}\small
\begin{tabular}{@{}lp{0.74\linewidth}@{}}
\toprule
Command & Effect \\
\midrule
$\shift$ & Halt if $t+h=1$. Otherwise replace $(t,w,h)$ by
$(h,p_R(w),p_L(w))$. \\
$\mpc(x)$ & Put $y=d+x$. If $h=I(t,y)$, return $(y,p(0,w),y,0)$.
Otherwise return $(0,w,y,0)$. \\
$\ax(c,x,y)$ & Put $a=d+x$ and $b=h+y$. Push eight axiom heads,
$U(d,h)$, and the saved top $t$. Set $t$ to the guarded separation
formula, or to truth when $c=a$ or $c=b$. Set $d=h=0$. \\
\bottomrule
\end{tabular}
\end{center}
The sixth head in this ordering is
\[
   \rho_6(a,b)=I\bigl(\ulcorner\exists x_a\,\forall x_a\,B\urcorner,b\bigr),
\]
where $B$ has code $b$ and the quantifiers use the numeric constructors.
The choice of arguments follows from the checked copy operations in
Section~\ref{sec:physical}. The other seven heads and the
generalization head keep their stated parameter dependence.

\subsection{Soundness and enumeration}

Interpret $y\in x$ in a model of $\PA$ as the statement that bit $y$ of
$x$ is set. The uniform axiom suppliers establish truth under this
interpretation, including pointed separation and the guarded
implicational schema. Modus ponens and generalization preserve
validity. The relation $y\in x$ implies $y<x$, so code $1$ is false
under every assignment in this interpretation. Soundness for an
actual finite derivation, followed by first-order completeness, gives
a $\PA$ refutation if code $1$ has a proof in the command calculus.

The clean store contains truth in the visible registers and in the
initial stack tail. The soundness invariant requires derivations for
the visible formulas and each entry in the relevant stack prefix.
An AX command supplies its heads with their derivations, an MP
command keeps the derived consequent or truth, and a Shift moves
derivations between the visible positions and stack. The only halt
test is $t+h=1$, which requires one visible formula to have code $1$.
The source scheduler preserves the invariant from the clean
candidate-$2$ store. It reads candidate codes by repeated unpairing,
consumes the command atoms in order, and proceeds to the next
candidate after a word finishes. The decoder and each candidate's
register computation terminate or take the designated halt branch.

Thus under $\Con(\PA)$ every candidate computation continues past
the halt test. Lemma~\ref{lem:simulation} proves blank-tape nonhalting.
This direction uses the external soundness theorem for actual finite
derivations in the command calculus.

\section{The exact arithmetic trace sentence}\label{sec:arith}

Encode a tape by natural numbers $l,r$. Their low bits are respectively
the cell just to the left of the head and the scanned cell. Higher bits
run away from the head. Missing high bits represent zero. For a
right-moving row $(q_0,b,w,R,q_1)$ use the arithmetic relation
\[
  q=q_0\land q'=q_1\land
  \exists u\,(r=2u+b\land l'=2l+w\land r'=u).
\]
For a left-moving row use
\begin{align*}
 q=q_0\land q'=q_1\land\exists u,v,e\,(&
    (e=0\lor e=1)\land l=2v+e\land r=2u+b\\
    &{}\land l'=v\land r'=2(2u+w)+e).
\end{align*}
Let $F(q,l,r,q',l',r')$ be the disjunction of these formulas over
the 418 rows. These formulas use the published state labels. Writing
sequence access in subscript notation, the halting sentence is
\begin{equation}\label{eq:halt}
 H:=\exists n,Q_*,L_*,R_*\,
 \left[
 \begin{aligned}
 &Q_0=L_0=R_0=0\land Q_n=215\\[-2pt]
 &{}\land\forall i<n\ F(Q_i,L_i,R_i,Q_{i+1},L_{i+1},R_{i+1})
 \end{aligned}
 \right].
\end{equation}
The three sequence variables range over arithmetic codes. Each access
in~\eqref{eq:halt} abbreviates the explicit arithmetic sequence
relation \texttt{nthDef}. In the standard model these relations
describe exactly the finite tape runs. In every model of $\PA$, the
exact formula is equivalent to the auxiliary formula used to compose
coded trace segments. For this equivalence, induction on the coded
step index excludes the six padding labels
$50,122,150,157,205,214$ from a trace starting at zero.
\src{M6PA/Model/ExactHaltingFormula.lean}{ExactHaltingFormula.lean}
proves the row correspondence and the two trace implications.

\section{From a coded PA refutation to halting}\label{sec:refutation}

Fix an arbitrary model $V$ of $\PA$. All objects in this section are
arithmetically coded in $V$. In particular, proof length, recursion
fuel, and trace length may be nonstandard.

\subsection{Refutation transport}

The arithmetized transport starts with a proof code for $\bot$ in
ordinary $\PA$. It translates the sequent proof to the presentation
with successor, and then interprets that presentation in the numeric
calculus of membership and equality. The interpretation uses the
finite ordinal domain with zero and ordinal successor, and treats
equality by the ordinal isomorphism relation. The axiom suppliers
realize the arithmetic defining equations and the induction scheme
by separation and ordinal induction. Logical inference is
transported by arithmetic operations on proof codes. This produces
an accepted certificate for code $1$ uniformly in $V$. The formal
transport is
\src{M6PA/Calibration/FoundationNativeRefutationTransport.lean}{FoundationNativeRefutationTransport.lean}.
It supplies the internal implication from standard $\PA$ provability
of $\bot$ to the existence of a certificate satisfying the exact
refutation-checker graph.

\subsection{Compiling an accepted certificate}\label{sec:compile}

A certificate is $p(f,i)$, with fuel $f$ and a tree code $i$.
A tree code pairs a rule tag with its payload. Payload tuples are
iterated Cantor pairs ending in zero. Tags $0,1$ are modus ponens and
generalization. Tags $2$ through $9$ are the eight axiom heads,
in order, and tag $10$ is the guarded separation instance.
Failed guards, malformed tuples and rejected premises reject the
certificate. The checker returns $a+1$ on accepting formula code
$a$, and returns zero on rejection. Thus the refutation-checker
graph tests output $2$.

The command compiler uses exact finite words to store, retrieve, and
discard formulas. In the following compositions, words execute from
left to right. The identity
\[
 \ax(c,x,y)\,\shift^{10}(t,w,d,h)=(U(d,h),w,0,t)
\]
follows by following the ten pushed items. The first ten Shift tests
see the gate and heads with residues $0$ or $2$ modulo $4$. The
following words give the required storage operations.
\begin{align*}
 \pophead&=\ax(0,0,0)\,\shift^{11},\\
 \push&=\ax(0,0,0)\,\shift^9\,\mpc(0),\\
 \discard&=\mpc(1)\,\ax(1,0,0)\,\shift^{10}\,\mpc(0).
\end{align*}
On a stack $p(a,w)$, the first word preserves $t$ and $w$ and returns
$h=a,d=0$. The second returns $(0,p(t,w),0,0)$ because the
generalization head cannot match the MP test. On $(t,w,0,0)$ the
third returns $(0,w,0,0)$, since its two MP tests reject.

To select head $j$ for $1\leq j\leq 7$, clear the top and use
\[
     \discard\,\ax(c,a,b)\,\shift^{j+1}\,\pophead^{9-j}.
\]
For the sixth head, use $\ax(0,a,b)$ to obtain $\rho_6(a,b)$.
For head eight with fourth parameter $u\leq a$, first set the
offset to $u$ after clearing the top and use $\ax(c,a-u,b)$.
Then use $\shift^9\pophead$. A nonzero offset $u$ is set by
$\mpc(u)$ on a zero top and zero caches. The zero offset requires
the empty word. The gate selector is
$\discard\,\ax(c,a,b)\,\pophead^{10}$. These words leave the
original stack tail and clear both caches.

Generalization of $a$ with label $i$ uses $\push$, sets offset $i$,
then executes $\shift\,\ax(i,0,0)\,\shift^{10}$. For MP with
premises $a$ and $I(a,b)$, compile the implication, execute $\push$,
compile $a$, and execute $\push$ again. The suffix is
\[
       \shift^2\,\mpc(b)\,\ax(b,0,0)\,\shift^{11}.
\]
The second Shift halts if $a=1$. Otherwise the MP succeeds, pushes
its zero, and the remaining commands return $b$ with the original
stack tail and zero caches. Each of these calculations uses
polynomial identities, total Cantor projections, and finite
command equations, so it holds in $V$.

For a formal induction in $V$, let $A(f,B)$ be the coded vector of
checker outputs for indices $0,\ldots,B$ at fuel $f$, and let
$W(f,B)$ be the vector of compiled words at the same indices.
The fuel-zero vectors contain rejections and empty words respectively.
A successor row reads the children from the previous vectors,
applies the rule just described, and appends one coded word as one
entry. Let $\mathcal R(w,a)$ mean that, for all $t,z$, running $w$
on $(t,z,0,0)$ halts or returns $(a,z,0,0)$. Full $\PA$ induction on
$f$ proves
\[
 \forall B,i,a\quad A(f,B)_i=a+1
                    \ \Longrightarrow\ \mathcal R(W(f,B)_i,a).
\]
At zero fuel the antecedent is false. At successor fuel an accepted
leaf uses its exact selector, a generalization uses the child
hypothesis followed by its suffix, and an MP uses both child
hypotheses followed by the MP suffix. Rejected rows cannot meet the
antecedent. This induction concerns an arithmetic predicate on coded
vectors and stores. It applies to all fuel values in $V$.

For a certificate accepting $1$, prefix the resulting word with
$\ax(0,0,0)$ and append a Shift. The prefix clears the two caches
from any initial store. If compilation has not already halted,
the last Shift sees $(t,h)=(1,0)$ and halts. This gives a
single word that halts from every store in $V$.

\subsection{Reaching the word and composing the trace}

The universal halting word begins with $\ax(0,0,0)$, so its first
atom has tag at least $2$. The word encoding gives a candidate index
at least $3$, after the actual startup at candidate $2$. The
scheduler's coded induction up to this index yields either an
earlier halt or a run of the universal word. This uses total-or-halt
for each earlier candidate and permits its resulting store to be
arbitrary.

The source-to-tape correspondence replaces every step in the coded
source trace by a positive-length coded tape block. It starts from
the full clean frame, concatenates these blocks, and prepends the
certified blank startup. The resulting coded trace ends at state
$215$. The row and trace correspondence in Section~\ref{sec:arith}
puts it in the exact formula $H$. Because $V$ was arbitrary,
first-order completeness gives the internal implication
\[
          \PA\vdash \Prov_{\PA}(\ulcorner\bot\urcorner)\to H,
          \qquad
          \PA\vdash N\to\Con(\PA).
\]
The compiled word, scheduler, and trace construction are checked
separately. The trace endpoint is
\src{M6PA/RefutationHalting.lean}{RefutationHalting.lean} and the final
consistency implication is in \src{M6PA/FinalEndpoint.lean}{FinalEndpoint.lean}.

\section{Two-sided independence}\label{sec:independence}

\begin{proof}[Proof of Theorem~\ref{thm:main}]
Suppose first that $\PA$ is consistent. The external search soundness
argument gives literal blank-tape nonhalting and hence $\N\models N$.
If $\PA$ proved $N$, the internal implication of
Section~\ref{sec:refutation} would give a $\PA$ proof of its own
syntactic consistency. G\"odel's second incompleteness
theorem~\cite{goedel} forbids this.

The standard natural-number structure satisfies the defining
equations of arithmetic and every induction instance. Soundness of
first-order derivations therefore yields external consistency and
truth in $\N$ of every $\PA$ theorem. In particular a $\PA$ proof of
$H$ would give $\N\models H$, contradicting $\N\models N$.
Thus $\PA\nvdash H$. Since classical $\PA$ proves $\neg N
\leftrightarrow H$, it also cannot prove $\neg N$. The standard model
discharges the consistency premise used in the first part.
\end{proof}

The Lean endpoint is
\src{M6PA/TwoSidedIndependence.lean}{TwoSidedIndependence.lean}.
\src{Solution.lean}{Solution.lean} transfers the exact formula and
provability predicate to the statements in
\src{Challenge.lean}{Challenge.lean}. The same Solution identifies
its tape semantics with the literal machine and transfers the
blank-tape nonhalting proof. The two statements concern the
particular machine defined by the table in Section~\ref{sec:machine}.

\begin{thebibliography}{9}

\bibitem{rado}
T. Rad\'o.
On non-computable functions.
\emph{Bell System Technical Journal} 41 (1962), 877--884.
\url{https://doi.org/10.1002/j.1538-7305.1962.tb00480.x}.

\bibitem{yedidia}
A. Yedidia and S. Aaronson.
A relatively small Turing machine whose behavior is independent of set theory.
\emph{Complex Systems} 25 (2016), 297--327.
\url{https://arxiv.org/abs/1605.04343}.

\bibitem{pa211}
J. Chang.
Nonhalting and PA-independence for a 211-state Turing machine.
\url{https://github.com/jamesyc/BusyBeaver/tree/1654314dd6892371ea38e488a7aa501d2aa55a0c}.

\bibitem{k230}
E. Glazer.
\emph{K230 Rocq verification}.
\url{https://github.com/elliotglazer/k230-rocq/tree/3142e4d}.

\bibitem{pa238}
\emph{PA238 Rocq development}.
\url{https://github.com/jeff87654/busy-beaver/tree/84049b5/PA/m238}.

\bibitem{proveit}
V. Reshetnikov.
\emph{ProveIt}, \texttt{DominatedCore.v}.
\url{https://github.com/VladimirReshetnikov/ProveIt/tree/6809609}.

\bibitem{goedel}
K. G\"odel.
\"Uber formal unentscheidbare S\"atze der Principia Mathematica und
verwandter Systeme I.
\emph{Monatshefte f\"ur Mathematik und Physik} 38 (1931), 173--198.
\url{https://doi.org/10.1007/BF01700692}.

\bibitem{foundation}
FormalizedFormalLogic.
\emph{Foundation: Formalized mathematical logic in Lean~4}.
\url{https://github.com/FormalizedFormalLogic/Foundation}.

\end{thebibliography}
\end{document}
